\chapter{CƠ SỞ LÝ THUYẾT}
\label{chap:theory}

[Chương này trình bày các khái niệm lý thuyết cốt lõi, cơ sở toán học và tổng quan các nghiên cứu liên quan đến đề tài.]

% ================= SECTION 2.1 =================
\section{Các khái niệm cơ bản}
\label{sec:concepts}

[Giải thích các định nghĩa, nguyên lý hoạt động nền tảng. Trích dẫn tài liệu học thuật liên quan như \cite{sample_conference, sample_standard}.]

% ================= SECTION 2.2 =================
\section{Mô hình và công thức toán học}
\label{sec:math}

[Trình bày các công thức và biến đổi toán học phục vụ cho hệ thống.]

Ví dụ công thức toán có đánh số:
\begin{equation}
    f(x) = \sum_{i=1}^{n} w_i \cdot x_i + b
    \label{eq:sample_linear}
\end{equation}
trong đó $w_i$ là hệ số trọng số, $x_i$ là biến đầu vào và $b$ là hệ số điều chỉnh (bias).

Biến đổi toán học nhiều dòng canh lề:
\begin{equation}
\begin{aligned}
    \mathcal{L}(\theta) &= -\frac{1}{N} \sum_{i=1}^{N} \left[ y_i \log(\hat{y}_i) + (1 - y_i) \log(1 - \hat{y}_i) \right] \\
    &= \frac{1}{N} \sum_{i=1}^{N} \ell_i(\theta)
\end{aligned}
\label{eq:sample_loss}
\end{equation}

Từ Phương trình~\eqref{eq:sample_linear} và Phương trình~\eqref{eq:sample_loss}, hệ thống tiến hành tối ưu hóa tham số trong quá trình huấn luyện.
